\documentclass[11pt]{article}

\usepackage{graphicx}
\usepackage{amsmath}
\usepackage{booktabs}
\usepackage[a4paper, margin=1.8cm]{geometry}
\usepackage{titling}
\usepackage{hyperref}
\usepackage{fontspec}
\usepackage{polyglossia}
\usepackage{titlesec}
\usepackage{caption}
\usepackage{float}
\usepackage{listings}
\usepackage{xcolor}

\DeclareCaptionLabelFormat{arabicfigure}{Նկար~\arabic{figure}}
\DeclareCaptionLabelFormat{arabictable}{Աղյուսակ~\arabic{table}}

\captionsetup[figure]{labelformat=arabicfigure,labelsep=period}
\captionsetup[table]{labelformat=arabictable,labelsep=period}

\setmainlanguage[numerals=arabic]{armenian}
\setotherlanguage{english}

\setmainfont{FreeSerif}
\newfontfamily\armenianfont{FreeSerif}

\setlength{\droptitle}{-2.5cm}

\title{\vspace{-1cm}Google Play Store հավելվածների բազմաչափ վիճակագրական վերլուծություն}
\author{Աննա Բեկլարյան}
\date{Հունիս, 2026}

\begin{document}

\renewcommand{\thefigure}{\arabic{figure}}
\renewcommand{\thetable}{\arabic{table}}
\renewcommand{\thepage}{\number\value{page}}

\renewcommand{\thesection}{\number\value{section}.}
\renewcommand{\thesubsection}{\number\value{section}.\number\value{subsection}.}

\titleformat{\section}
  {\normalfont\Large\bfseries}
  {\thesection}
  {0.3em}
  {}

\titleformat{\subsection}
  {\normalfont\large\bfseries}
  {\thesubsection}
  {0.3em}
  {}

\maketitle
\vspace{-1cm}

\section{Ներածություն և տվյալների բազայի նկարագրություն}
Այս նախագծի համար ընտրվել է \textit{Google Play Store Apps Dataset} տվյալների բազան\footnote{\url{https://www.kaggle.com/datasets/lava18/google-play-store-apps}}, որը պարունակում է տեղեկություններ \textbf{ավելի քան 10,000 Android հավելվածների} մասին, որոնք հասանելի են Google Play Store հարթակում։ Տվյալների բազան ներառում է հավելվածների վերաբերյալ տարբեր բնութագրեր, ինչպիսիք են օգտատերերի գնահատականները, կարծիքների քանակը, ներբեռնումների քանակը, հավելվածի չափը, գինը և այլն։

Վիճակագրական վերլուծության նպատակով ընտրվել են հետևյալ քանակական փոփոխականները.

\begin{itemize}
\item \textbf{Rating} -- հավելվածի օգտատերերի միջին գնահատականը։
\item \textbf{Reviews} -- օգտատերերի կարծիքների քանակը։
\item \textbf{Size} -- հավելվածի չափը մեգաբայթերով (MB)։
\item \textbf{Installs} -- հավելվածի ներբեռնումների քանակը։
\item \textbf{Price} -- հավելվածի գինը ԱՄՆ դոլարով։
\end{itemize}

Այս տվյալների բազան հետաքրքիր է վիճակագրական վերլուծության համար, քանի որ բջջային հավելվածները լայնորեն օգտագործվում են առօրյա կյանքում, իսկ դրանց հաջողությունը կարող է կախված լինել տարբեր չափելի գործոններից՝ գնահատականներից, կարծիքների քանակից, ներբեռնումներից, չափից և գնից։ Տվյալների բազան հնարավորություն է տալիս ուսումնասիրել այդ գործոնների միջև առկա կապերը և հասկանալ, թե որոնք են առավելապես կապված հավելվածների հանրաճանաչության և օգտատերերի բավարարվածության հետ։ Բացի այդ, ընտրված քանակական փոփոխականները հարմար են ռեգրեսիոն վերլուծության և հիմնական բաղադրիչների վերլուծության (Principal Component Analysis, PCA) կիրառման համար։

\section{Տվյալների կառուցվածքի նախնական ուսումնասիրություն}

\subsection{Կետային ամպերի մատրից}

Նախ ընտրված քանակական փոփոխականների համար կառուցվել է երկչափ կետային ամպերի մատրից։ Քանի որ \textbf{Reviews} և \textbf{Installs} փոփոխականները ունեն շատ մեծ արժեքներ, գրաֆիկում օգտագործվել են դրանց լոգարիթմացված տարբերակները՝ \textbf{log\_Reviews} և \textbf{log\_Installs}։ Սա թույլ է տալիս ավելի հստակ պատկերել տվյալների կառուցվածքը։

Նկար~\ref{fig:scatter_matrix}-ում երևում է, որ ամենահստակ դրական կապը առկա է \textbf{log\_Reviews} և \textbf{log\_Installs} փոփոխականների միջև։ Սա բնական է, քանի որ շատ ներբեռնումներ ունեցող հավելվածները սովորաբար ունեն նաև ավելի շատ օգտատերերի կարծիքներ։ \textbf{Rating} փոփոխականի կապը մյուս փոփոխականների հետ ավելի թույլ է, իսկ \textbf{Size} փոփոխականի դեպքում հստակ գծային կապ չի նկատվում։

\begin{figure}[H]
\centering
\includegraphics[width=0.52\textwidth]{scatter_matrix_googleplay.png}
\caption{Ընտրված փոփոխականների կետային ամպերի մատրից}
\label{fig:scatter_matrix}
\end{figure}

\subsection{Բազմաչափ նորմալության ենթադրության գնահատում}

Ստացված կետային ամպերի և անկյունագծի վրա պատկերված բաշխումների հիման վրա կարելի է գնահատել, թե որքանով է հիմնավորված բազմաչափ նորմալ բաշխման ենթադրությունը։

Նկատվում է, որ որոշ փոփոխականների բաշխումներ շեղված են, իսկ կետային ամպերը հիմնականում չունեն էլիպսաձև տեսք, որը բնորոշ է բազմաչափ նորմալ բաշխմանը։ Հատկապես \textbf{Reviews}, \textbf{Installs} և \textbf{Size} փոփոխականների դեպքում նկատվում են ասիմետրիկ բաշխումներ և ոչ նորմալ կառուցվածքներ։

Հետևաբար կարելի է եզրակացնել, որ տվյալները ընդհանուր առմամբ չեն բավարարում բազմաչափ նորմալ բաշխման ենթադրությանը։

\section{Ռեգրեսիոն վերլուծություն}

Որպես ռեգրեսիայի բացատրվող փոփոխական ընտրվել է \textbf{Rating} փոփոխականը, քանի որ այն ներկայացնում է հավելվածի օգտատերերի միջին գնահատականը և բնութագրում է օգտատերերի բավարարվածության մակարդակը։ Որպես բացատրող փոփոխականներ օգտագործվել են \textbf{Reviews}, \textbf{Size}, \textbf{Installs} և \textbf{Price} փոփոխականները։ Քանի որ \textbf{Reviews}, \textbf{Installs} և \textbf{Price} փոփոխականներն ունեն խիստ շեղված բաշխումներ, դրանց նկատմամբ կիրառվել է լոգարիթմական ձևափոխություն։

Կառուցվել է հետևյալ բազմակի գծային ռեգրեսիայի մոդելը.
\[
\textit{Rating}
=
\beta_0
+
\beta_1 \log(\textit{Reviews})
+
\beta_2 \textit{Size}
+
\beta_3 \log(\textit{Installs})
+
\beta_4 \log(\textit{Price})
+
\varepsilon
\]

Մոդելի պարամետրերը գնահատվել են նվազագույն քառակուսիների մեթոդով (Ordinary Least Squares, OLS).
\[
\hat{\beta}_{OLS} = (X^TX)^{-1}X^TY
\]
որտեղ \(X\)-ը բացատրող փոփոխականների մատրիցն է, \(Y\)-ը բացատրվող փոփոխականի վեկտորը, իսկ \(\hat{\beta}_{OLS}\)-ը գնահատված ռեգրեսիոն գործակիցների վեկտորը։

\begin{figure}[H]
\centering
\begin{minipage}{0.48\textwidth}
    \centering
    \includegraphics[width=\textwidth]{regression_results.png}
    \caption*{(ա) Սկզբնական մոդել}
\end{minipage}
\hfill
\begin{minipage}{0.51\textwidth}
    \centering
    \includegraphics[width=\textwidth]{regression_selected.png}
    \caption*{(բ) Փոփոխականների ընտրությունից հետո}
\end{minipage}
\caption{OLS ռեգրեսիայի արդյունքները փոփոխականների ընտրությունից առաջ և հետո}
\label{fig:regression_comparison}
\end{figure}

Սկզբնական մոդելի գնահատված հավասարումն ունի հետևյալ տեսքը.
\[
\widehat{\textit{Rating}}
=
4.6015
+
0.1653 \log(\textit{Reviews})
-
0.0002 \textit{Size}
-
0.1453 \log(\textit{Installs})
-
0.0420 \log(\textit{Price})
\]

Մոդելի \(R^2\) գործակիցը հավասար է \(0.116\), ինչը նշանակում է, որ ընտրված բացատրող փոփոխականները միայն մասնակիորեն են բացատրում հավելվածների գնահատականների տարբերությունները։ Այսինքն՝ \textbf{Rating}-ի վրա կարող են ազդել նաև այլ գործոններ, որոնք ներառված չեն տվյալ մոդելում։

Գործակիցների վերլուծությունը ցույց է տալիս, որ \textbf{log(Reviews)} փոփոխականը դրական կապ ունի \textbf{Rating}-ի հետ։ Միևնույն ժամանակ, \textbf{log(Installs)} և \textbf{log(Price)} փոփոխականների գործակիցները բացասական են։ \textbf{Size} փոփոխականի գործակիցը շատ փոքր է և վիճակագրորեն նշանակալի չէ (\(p=0.385\)), հետևաբար հավելվածի չափի ազդեցությունը գնահատականի վրա տվյալների հիման վրա չի հաստատվում։

Ընդհանուր առմամբ, մոդելը վիճակագրորեն նշանակալի է (\(F=253.6\), \(p<0.001\)), ինչը ցույց է տալիս, որ ընտրված բացատրող փոփոխականներից առնվազն մեկը վիճակագրորեն կապված է \textbf{Rating} փոփոխականի հետ։

Քանի որ \textbf{Size} փոփոխականը վիճակագրորեն նշանակալի չէր, այն հեռացվել է մոդելից։ Վերջնական մոդելի գնահատված հավասարումն է.
\[
\widehat{\textit{Rating}}
=
4.5987
+
0.1644 \log(\textit{Reviews})
-
0.1449 \log(\textit{Installs})
-
0.0420 \log(\textit{Price})
\]

Վերջնական մոդելի \(R^2\) գործակիցը կրկին հավասար է \(0.116\), ինչը ցույց է տալիս, որ \textbf{Size} փոփոխականի հեռացումը չի նվազեցրել մոդելի բացատրական կարողությունը։ Վերջնական մոդելում մնացած բոլոր փոփոխականները վիճակագրորեն նշանակալի են, ուստի հետագա մեկնաբանության համար ընտրվել է այս մոդելը։

\section{Հիմնական բաղադրիչների վերլուծություն}

Հիմնական բաղադրիչների վերլուծությունը (PCA) կատարվել է ընտրված քանակական փոփոխականների վրա՝ \textbf{Rating}, \textbf{log(Reviews)}, \textbf{Size}, \textbf{log(Installs)} և \textbf{log(Price)}։ Քանի որ փոփոխականները տարբեր մասշտաբներ ունեն, նախ դրանք ստանդարտացվել են։

PCA-ի նպատակն է մի քանի փոխկապակցված փոփոխականները փոխարինել նոր փոփոխականներով՝ հիմնական բաղադրիչներով, որոնք բացատրում են տվյալների տատանումների առավել մեծ մասը։

\begin{table}[H]
\centering
\caption{Հիմնական բաղադրիչների բացատրած տատանումները}
\label{tab:pca_variance}
\begin{tabular}{ccc}
\toprule
Բաղադրիչ & Տատանում (\%) & Կուտակային (\%) \\
\midrule
PC1 & 45.1 & 45.1 \\
PC2 & 20.8 & 65.8 \\
PC3 & 18.3 & 84.2 \\
PC4 & 15.1 & 99.2 \\
PC5 & 0.8 & 100.0 \\
\bottomrule
\end{tabular}
\end{table}

Աղյուսակ~\ref{tab:pca_variance}-ից երևում է, որ առաջին երկու հիմնական բաղադրիչները միասին բացատրում են ընդհանուր տատանումների մոտ \textbf{65.8\%}-ը, իսկ առաջին երեք բաղադրիչները՝ մոտ \textbf{84.2\%}-ը։

\begin{figure}[H]
\centering
\includegraphics[width=0.5\textwidth]{pca_scree_plot.png}
\caption{PCA-ի Scree plot}
\label{fig:pca_scree}
\end{figure}

Նկար~\ref{fig:pca_scree}-ը ցույց է տալիս, որ առաջին հիմնական բաղադրիչն ունի ամենամեծ բացատրական ուժը։ Հինգերորդ բաղադրիչի ներդրումը շատ փոքր է, հետևաբար տվյալների հիմնական կառուցվածքը կարելի է նկարագրել առաջին մի քանի բաղադրիչների միջոցով։

\begin{figure}[H]
\centering
\includegraphics[width=0.5\textwidth]{pca_correlation_circle.png}
\caption{PCA-ի փոփոխականների կորելացիաների շրջան}
\label{fig:pca_circle}
\end{figure}

\begin{table}[H]
\centering
\caption{Առաջին երկու հիմնական բաղադրիչների բեռնվածքները}
\label{tab:pca_loadings}
\begin{tabular}{lcc}
\toprule
Փոփոխական & PC1 & PC2 \\
\midrule
Rating & -0.158 & 0.672 \\
log(Reviews) & -0.636 & 0.022 \\
Size & -0.374 & 0.215 \\
log(Installs) & -0.628 & -0.106 \\
log(Price) & 0.191 & 0.700 \\
\bottomrule
\end{tabular}
\end{table}

Նկար~\ref{fig:pca_circle}-ում ներկայացված է փոփոխականների կորելացիաների շրջանը, իսկ Աղյուսակ~\ref{tab:pca_loadings}-ում՝ առաջին երկու բաղադրիչների բեռնվածքները։ \textbf{log(Reviews)} և \textbf{log(Installs)} փոփոխականները գրեթե նույն ուղղությամբ են, ինչը վկայում է դրանց միջև ուժեղ դրական կապի մասին։ Սա բնական է, քանի որ շատ ներբեռնումներ ունեցող հավելվածները սովորաբար ունեն նաև շատ կարծիքներ։

Առաջին հիմնական բաղադրիչը հիմնականում կապված է \textbf{log(Reviews)} և \textbf{log(Installs)} փոփոխականների հետ, քանի որ դրանց բեռնվածքները PC1-ում մեծ են բացարձակ արժեքով։ Հետևաբար PC1-ը կարելի է մեկնաբանել որպես հավելվածների հանրաճանաչության ուղղություն։

Երկրորդ հիմնական բաղադրիչը հիմնականում կապված է \textbf{Rating} և \textbf{log(Price)} փոփոխականների հետ։ Սա նշանակում է, որ PC2-ը ավելի շատ արտահայտում է հավելվածի գնահատականի և գնի հետ կապված տարբերությունները։

\begin{figure}[H]
\centering
\includegraphics[width=0.6\textwidth]{pca_observations_projection.png}
\caption{Տվյալների պրոյեկցիան առաջին երկու հիմնական բաղադրիչների վրա}
\label{fig:pca_projection}
\end{figure}

Նկար~\ref{fig:pca_projection}-ում ներկայացված է հավելվածների պրոյեկցիան առաջին երկու հիմնական բաղադրիչների վրա։ Յուրաքանչյուր կետ համապատասխանում է մեկ հավելվածի։

Գրաֆիկից երևում է, որ հավելվածների մեծ մասը կենտրոնացած է գրաֆիկի կենտրոնական հատվածում։ Սա նշանակում է, որ ընտրված փոփոխականների տեսանկյունից հավելվածների մեծ մասը իրարից շատ կտրուկ չի տարբերվում։ Միաժամանակ կան կենտրոնից հեռու գտնվող կետեր, որոնք կարելի է դիտարկել որպես առանձնացող դիտարկումներ (outliers)։

Տարածվածությունն ավելի մեծ է առաջին հիմնական բաղադրիչի ուղղությամբ, ինչը բնական է, քանի որ PC1-ը բացատրում է տվյալների ընդհանուր տատանումների ավելի մեծ մասը, քան PC2-ը։



\section{Կորելացիոն վերլուծություն}

Որպես լրացուցիչ վիճակագրական մեթոդ իրականացվել է կորելացիոն վերլուծություն։ Դրա նպատակն է գնահատել ընտրված քանակական փոփոխականների միջև առկա գծային կապերի ուժգնությունն ու ուղղությունը։

\begin{figure}[H]
\centering
\includegraphics[width=0.45\textwidth]{correlation_heatmap.png}
\caption{Ընտրված փոփոխականների կորելացիոն մատրից}
\label{fig:correlation_heatmap}
\end{figure}

Նկար~\ref{fig:correlation_heatmap}-ում ներկայացված է ընտրված փոփոխականների կորելացիոն մատրիցը։ Ամենաուժեղ դրական կապը դիտվում է \textbf{log(Reviews)} և \textbf{log(Installs)} փոփոխականների միջև, որոնց կորելացիայի գործակիցը հավասար է \textbf{0.95}-ի։ Սա բնական արդյունք է, քանի որ մեծ թվով ներբեռնումներ ունեցող հավելվածները սովորաբար ունեն նաև մեծ թվով օգտատերերի կարծիքներ։

\textbf{Size} փոփոխականը ցուցաբերում է չափավոր դրական կապ \textbf{log(Reviews)} (\(r=0.38\)) և \textbf{log(Installs)} (\(r=0.33\)) փոփոխականների հետ, ինչը նշանակում է, որ ավելի մեծ չափ ունեցող հավելվածները որոշ չափով հակված են ունենալ ավելի շատ ներբեռնումներ և կարծիքներ։

\textbf{Rating} փոփոխականի և մյուս փոփոխականների միջև կապերը համեմատաբար թույլ են։ Ամենամեծ կորելացիան դիտվում է \textbf{log(Reviews)} փոփոխականի հետ (\(r=0.18\)), սակայն այն նույնպես թույլ է։ Սա ցույց է տալիս, որ հավելվածի գնահատականը միայն փոքր չափով է կապված ընտրված քանակական բնութագրերի հետ։

Բացի այդ, \textbf{log(Price)} փոփոխականը ցուցաբերում է թույլ բացասական կապ \textbf{log(Reviews)} (\(r=-0.14\)) և \textbf{log(Installs)} (\(r=-0.22\)) փոփոխականների հետ։ Այս արդյունքը կարող է վկայել այն մասին, որ ավելի թանկ հավելվածները սովորաբար ունեն ավելի քիչ ներբեռնումներ և օգտատերերի կարծիքներ։

Ընդհանուր առմամբ, կորելացիոն վերլուծությունը հաստատում է, որ տվյալների կառուցվածքում առավել կարևոր կապը գոյություն ունի \textbf{Reviews} և \textbf{Installs} փոփոխականների միջև, մինչդեռ \textbf{Rating} փոփոխականը համեմատաբար թույլ է կապված մյուս ընտրված փոփոխականների հետ։


\section{Եզրակացություն}

Այս աշխատանքում իրականացվել է Google Play Store հավելվածների բազմաչափ վիճակագրական վերլուծություն։

Երկչափ կետային ամպերի ուսումնասիրությունը ցույց տվեց, որ \textbf{log(Reviews)} և \textbf{log(Installs)} փոփոխականների միջև առկա է ուժեղ դրական կապ։ Միևնույն ժամանակ, փոփոխականների բաշխումները և կետային ամպերը ցույց տվեցին, որ տվյալները ընդհանուր առմամբ չեն համապատասխանում բազմաչափ Գաուսյան բաշխմանը։

Ռեգրեսիոն վերլուծության համար որպես բացատրվող փոփոխական ընտրվել է \textbf{Rating}-ը։ Ստացված մոդելը ցույց տվեց, որ \textbf{log(Reviews)}, \textbf{log(Installs)} և \textbf{log(Price)} փոփոխականները վիճակագրորեն նշանակալի են, մինչդեռ \textbf{Size} փոփոխականի ազդեցությունը նշանակալի չէ։ Միաժամանակ \(R^2=0.116\) արժեքը ցույց է տալիս, որ գնահատականի վրա ազդում են նաև այլ գործոններ, որոնք տվյալ մոդելում ներառված չեն։

PCA-ի արդյունքները ցույց տվեցին, որ առաջին երկու հիմնական բաղադրիչները բացատրում են տվյալների ընդհանուր տատանումների մոտ \textbf{65.8\%}-ը։ Առաջին բաղադրիչը հիմնականում կապված է հավելվածների հանրաճանաչության հետ, իսկ երկրորդը՝ գնահատականի և գնի հետ։

Ընդհանուր առմամբ, վերլուծությունը ցույց է տալիս, որ հավելվածների ներբեռնումների և կարծիքների քանակը տվյալների կառուցվածքը ձևավորող հիմնական գործոններից են, իսկ հավելվածի գնահատականը միայն մասամբ է բացատրվում ընտրված քանակական փոփոխականներով։

\appendix

\section{Հավելված. Python կոդ}

\lstinputlisting[language=Python]{analysis.py}

\section{Կոդի հասանելիություն}

Աշխատանքում օգտագործված ամբողջ Python կոդը, Jupyter Notebook ֆայլերը և ստացված պատկերները հասանելի են հետևյալ GitHub պահոցում.

\url{https://github.com/AnnaBeklaryan/google-play-mv-stats}

\end{document}

